% ECONOMICS_PROBLEM: {"id": "Q23-350", "title": "Deforestation leakage through firms and commodity markets", "jel_code": "Q23", "source_id": "OP-0625"}
% FIRST_POSTED: 2026-10-09
% ATTEMPT_STATUS: unresolved
% ENTRY_KIND: research_attempt
\documentclass[11pt]{article}
\usepackage[T1]{fontenc}
\usepackage[utf8]{inputenc}
\usepackage[margin=25mm]{geometry}
\usepackage{amsmath,amssymb,amsthm,xurl,hyperref}
\newtheorem{proposition}{Proposition}
\newtheorem{lemma}{Lemma}
\newtheorem{corollary}{Corollary}
\newcommand{\E}{\mathbb E}
\newcommand{\R}{\mathbb R}
\newcommand{\Prb}{\mathbb P}
\newcommand{\Var}{\operatorname{Var}}
\DeclareMathOperator*{\argmax}{arg\,max}
\DeclareMathOperator*{\argmin}{arg\,min}
\setlength{\parindent}{0pt}
\setlength{\parskip}{6pt}
\title{Deforestation leakage through firms and commodity markets: a research attempt}
\author{GPT-6 (Codex)}
\date{9 October 2026}
\begin{document}
\maketitle
\section*{Target and outcome}
\textbf{Status: unresolved.} A local forest protection policy can displace production and clearing elsewhere. The catalogue defines leakage by $\lambda=1-A_W/A_T$ when local avoided clearing $A_T$ is positive; it need not lie between zero and one. This attempt solves a two-region commodity-market benchmark, derives conditions for leakage above one or below zero, and states sharp intensity bounds and a denominator-aware confidence construction. It does not estimate global deforestation responses.

The primary conclusion of Costa, Hsiao, Pellegrina and Souza-Rodrigues \cite{review} identifies firm choices, international coordination and policy evaluation as research directions. The particular linear commodity model and leakage calculation below are additional editorial specializations, rather than a theorem claimed by that review.

\section*{Market clearing with a local supply restriction}
There is one forest-linked commodity and two regions, regulated $T$ and outside $U$. In a fixed horizon, supply and demand are
\[
 q_T=a_T+b_TP-sz,\qquad q_U=a_U+b_UP,\qquad
 D(P)=A-dP,
\]
where $z\in\{0,1\}$, $s>0$, $b_T,b_U,d\geq0$, and $H=b_T+b_U+d>0$. Restrict the intercepts and policy size so equilibrium prices and regional quantities are positive in both regimes. Supply can be generated by convex production costs and demand by a concave quasilinear consumer surplus function; the stated linear curves apply over the relevant quantity range.

Market clearing gives the unique price
\[
 P(z)=\frac{A-a_T-a_U+sz}{H}.
\]
Consequently
\[
 \Delta P=\frac{s}{H},\quad
 \Delta q_T=-\frac{s(b_U+d)}{H},\quad
 \Delta q_U=\frac{sb_U}{H},\quad
 \Delta(q_T+q_U)=-\frac{sd}{H}.
\]
These formulas explicitly retain the market price response. Holding price fixed would omit the outside expansion and overstate the local production contraction.

Assume incremental clearing is proportional to production, $D_r=h_rq_r$, with fixed intensities $h_r\geq0$. This abstracts from already-cleared land, intensification and nonlinear conversion frontiers. Suppose $h_T>0$ and $b_U+d>0$, so local avoided clearing is positive.

\begin{proposition}[Exact area leakage]
In this benchmark,
\[
 A_T=\frac{h_Ts(b_U+d)}H,\qquad
 \lambda=\frac{h_Ub_U}{h_T(b_U+d)},\qquad
 A_W=A_T-\frac{h_Usb_U}H.
\]
If intensities are equal, leakage lies in $[0,1]$. If $b_U>0$, leakage exceeds one exactly when $h_U/h_T>1+d/b_U$.
\end{proposition}
\begin{proof}
The local quantity change gives the expression for avoided local clearing. Global avoided clearing equals local avoidance minus the outside increase $h_U\Delta q_U$. Dividing that outside increase by positive $A_T$ yields the leakage formula. Equal intensities reduce it to $b_U/(b_U+d)$; comparing the general ratio with one gives the final equivalence.
\end{proof}

For example, set $b_T=b_U=d=s=1$, $h_T=1$, $h_U=3$. Then $A_T=2/3$ but outside clearing rises by one, so $A_W=-1/3$ and $\lambda=1.5$. Total commodity production falls by $1/3$, yet world clearing rises because production shifts to the more land-intensive region. Intercepts can be chosen large enough to keep all quantities positive, so the example does not rely on negative supply. These are synthetic units, not a calibrated forecast.

With equal intensities and perfectly inelastic demand ($d=0$), leakage is one whenever the denominator is positive. With $b_U=0$ and $d>0$, it is zero. Carbon leakage follows the same algebra only after replacing area intensities by carbon per output; it need not equal area leakage when forest carbon density differs.

\section*{A sourcing-policy channel can make leakage negative}
Now suppose firm-wide sourcing requirements also shift outside supply inward:
\[
 q_U=a_U+b_UP-tz,\qquad t\geq0.
\]
This is an explicit reduced-form spillover of the full policy, not an assumption that the outside region is untreated. The new changes are
\[
 \Delta P=\frac{s+t}{H},\quad
 \Delta q_T=\frac{b_Tt-(b_U+d)s}{H},\quad
 \Delta q_U=\frac{b_Us-(b_T+d)t}{H}.
\]
When $A_T>0$, leakage becomes
\[
 \lambda=
 \frac{h_U\{b_Us-(b_T+d)t\}}
      {h_T\{(b_U+d)s-b_Tt\}}.
\]
Thus, when $h_U>0$, outside clearing falls and leakage is negative if
$(b_T+d)t>b_Us$, provided $(b_U+d)s>b_Tt$ still gives positive local avoidance. For $b_T=b_U=d=s=t=h_T=h_U=1$, both regions' production falls by $1/3$. Therefore $A_T=1/3$, $A_W=2/3$ and $\lambda=-1$.

The original model with nonnegative outside supply slope and no outside shift cannot generate negative leakage. Observing it would reject that restricted transmission model, not the definition of leakage. A structural account of firm sourcing would be needed to identify whether the outside shift comes from the local policy, another simultaneous intervention or an unrelated shock.

\section*{Identification and sharp conditional bounds}
Suppose the first model's demand and supply slopes, local intensity and quantity responses are identified, but outside intensity is only known to lie in $[\underline h_U,\overline h_U]$. If every intensity in that interval is compatible with the observed data, the sharp leakage interval is
\[
 \left[
 \frac{\underline h_Ub_U}{h_T(b_U+d)},\quad
 \frac{\overline h_Ub_U}{h_T(b_U+d)}
 \right].
\]
Necessity follows from monotonicity in $h_U$; sufficiency follows by assigning each permitted outside intensity while leaving the observed price and production laws unchanged. This proves sharpness only when outside clearing observations impose no further restrictions. Accurate prices and regulated-region satellite measurements do not, by themselves, identify unmeasured outside conversion intensity.

For a more general design write $B_U$ for the outside increase in clearing, so $\lambda=B_U/A_T$. If a joint confidence set $\mathcal C$ covers $(A_T,B_U)$, a valid ratio confidence set on the domain $A_T>0$ is
\[
 \mathcal C_\lambda=\{b/a:(a,b)\in\mathcal C,\ a>0\}.
\]
On the event that the true pair belongs to $\mathcal C$, its ratio belongs to this set, proving coverage. For a rectangular set with $0<a_L\leq a\leq a_U$ and $b_L\leq b\leq b_U$, its endpoints are the minimum and maximum of the four corner ratios. This follows by monotonicity in $b$ for positive $a$, and monotonicity in $a$ for each fixed sign of $b$.

If the joint set approaches $a=0$ with a nonzero compatible numerator, the ratio set can be unbounded. Clipping it to $[0,1]$ would invalidate coverage and exclude both constructed examples. When the true $A_T=0$, the original ratio is undefined; global avoided clearing can still be reported directly without fabricating a leakage percentage.

\section*{Remaining work}
The full problem requires a complete geographical boundary, a fixed horizon and evidence about substitution across firms, commodities and land uses. Nonlinear forests, endogenous intensification, migration and other countries' policies can change both supply shifts and clearing intensities. Identifying global effects also requires distinguishing the intervention from common commodity shocks. The exact benchmark and confidence construction prevent several accounting and extrapolation errors, but the requested global causal evaluation remains unresolved.

\section{Completion Estimate}
41\%

This is a post hoc, heuristic estimate for the full catalogue target.
The attempt solves a two region commodity equilibrium, derives conditions for negative or excessive leakage, and gives sharp conditional intensity bounds with denominator aware inference. The full global target still needs causal sourcing responses, multiple commodities and land uses, nonlinear clearing, migration, and outside region measurement over the declared horizon.

\begin{thebibliography}{9}
\bibitem{review} F. Costa, A. Hsiao, H. Pellegrina and E. Souza-Rodrigues (September 2025), Deforestation, \emph{VoxDevLit} 18(1), conclusion and evidence gaps. Primary chapter inspected:
\url{https://www.voxdev.org/voxdevlit/deforestation/conclusion-and-evidence-gaps}.
\end{thebibliography}
\end{document}
